%% ====================================================================
\section{Relation to Prior Work}
\label{sec:prior}
%% ====================================================================

\subsection{Odlyzko's Numerical Verification}
\label{sec:odlyzko}

Andrew Odlyzko's seminal computations~\cite{odlyzko1989,odlyzko2001} focused on
computing specific high zeros of~$\zeta$---around the $10^{20}$th and $10^{22}$nd
zeros---and large blocks of consecutive zeros near these heights, providing crucial
statistical data on zero spacing and GUE pair correlation. The systematic verification
that the first $10^{13}$ nontrivial zeros of~$\zeta$ lie on the critical line was carried
out by Gourdon~\cite{gourdon2004}. Both Odlyzko's and Gourdon's methods use the
Odlyzko--Sch\"onhage algorithm, which evaluates $\zeta(1/2 + it)$ at many points
simultaneously using the Riemann--Siegel formula and FFT-based polynomial evaluation,
achieving $O(T^{1/2+\eps})$ complexity for all zeros up to height~$T$. These remain the
gold standard for numerical verification of~RH.

Our work does not compete with Odlyzko's in terms of the number of zeros verified.
Instead, it provides:
\begin{itemize}
  \item An alternative computational framework based on quantum physics.
  \item A cross-check using an entirely different mathematical representation (eta
        function and $Z$-function via DQPT order parameters rather than direct zeta
        evaluation).
  \item Visualization of the $\beta$-dependence (the 2D phase diagram), which is not part
        of the standard numerical verification workflow.
\end{itemize}


\subsection{Berry--Keating Conjecture}
\label{sec:berry-keating}

Berry and Keating~\cite{berrykeating1999} conjectured that the Hilbert--P\'olya operator
for~RH is related to the quantization of the classical Hamiltonian $H = xp$. Specifically,
they proposed that the operator $(xp + px)/2$, suitably regularized, has eigenvalues
related to the nontrivial zeta zeros. This connects~RH to quantum chaos: the zeta zeros
would be energy levels of a chaotic quantum system, and their statistics (Montgomery's
pair correlation conjecture~\cite{montgomery1973}, confirmed numerically by Odlyzko)
match those of random matrix theory (GUE ensemble).

The Wei et~al.\ construction~\cite{wei2025} is distinct from Berry--Keating in that it
does not seek an operator whose \emph{spectrum} encodes the zeros, but rather constructs
systems whose \emph{dynamical phase transitions} (temporal non-analyticities) correspond
to zeros. The spectrum of $H_0 = \sum\ln(n)\,|n\rangle\langle n|$ is logarithmic by
design, and the zeros emerge in the time domain rather than the energy domain.


\subsection{Connes' Operator-Algebraic Approach}
\label{sec:connes}

Alain Connes~\cite{connes1999} formulated a trace formula relating the zeros of~$\zeta$
to the spectrum of a certain operator on an adelic space, providing a noncommutative
geometry framework for~RH. This approach has led to deep structural insights but has not
yet yielded a proof of~RH. Our computational work is complementary: it provides numerical
data that any future proof (whether through Connes' framework or otherwise) must be
consistent with.


\subsection{Guth--Maynard 2024 Breakthrough}
\label{sec:guth-maynard}

In 2024, Larry Guth and James Maynard~\cite{guthmaynard2024} proved improved large-value
estimates for Dirichlet polynomials, which imply that at most $O(T^{3/5+\eps})$ zeros of
$\zeta(s)$ with $0 < \im(s) \le T$ can lie off the critical line, improving the previous
$O(T^{3/4+\eps})$ bound (Ingham~\cite{ingham1940}; refined by Huxley and others). Note
that the qualitative conclusion---that $100\%$ of zeros lie on the critical line in a
density sense (i.e., $N_{\mathrm{off}}(T)/N(T) \to 0$)---already followed from the prior
$O(T^{3/4+\eps})$ bound since $T^{3/4}/N(T) \to 0$. Guth--Maynard's contribution is the
improved quantitative exponent, which represents the first improvement to the zero density
estimate in the critical strip in over $80$~years.

Our numerical investigation probes the complementary question: among the specific zeros
in our scanned range, are there any off the critical line? Guth--Maynard tells us that
at most $T^{3/5+\eps}$ could be off-line up to height~$T$; our computation would detect
any such zeros within $[0, T_{\max}]$.


\subsection{How Our Implementation Differs}
\label{sec:novelty}

Our implementation is, to our knowledge, the first public classical simulation of the
Wei et~al.\ DQPT framework~\cite{wei2025}. It differs from all the above in:
\begin{enumerate}
  \item \textbf{Physical framing.} We compute quantum-mechanical observables (coherence
        functions, Loschmidt echoes, rate functions) rather than directly evaluating
        $\zeta$ or $L$-functions.
  \item \textbf{2D scanning.} We scan the full $(\beta, t)$ plane, not just the critical
        line.
  \item \textbf{DQPT characterization.} We identify zeros through their physical
        manifestation as phase transitions (singularities in the rate function), not
        merely as zero-crossings.
  \item \textbf{Dual systems.} We implement both System~1 (eta-based) and System~2
        ($Z$-function-based), providing internal cross-validation.
\end{enumerate}
